% !TEX root = ../linnik399.tex
\section{Introduction}\label{sec:intro}

\subsection{The result}
For an integer $q\ge2$ and an integer $a$ with $(a,q)=1$, let $P(a,q)$ be the least prime
congruent to $a$ modulo $q$. Dirichlet's theorem \cite{Dirichlet1889beweis} ensures that this
prime exists. Linnik \cite{Linnik1944least,Linnik1944deuring} proved that there are absolute
constants $C$ and $L$ such that
\begin{equation}\label{eq:linnik}
  P(a,q)\le Cq^L
\end{equation}
for every reduced residue class. An exponent $L$ for which such a constant $C$ exists is
called \emph{admissible}; Linnik's constant is the infimum of the admissible exponents.

\begin{theorem}[Least prime in a reduced residue class]\label{thm:main}
There is an absolute constant $q_0\ge2$ such that, for every integer $q\ge q_0$ and every
integer $a$ with $(a,q)=1$, there is a prime $p$ satisfying
\[
  p\equiv a\pmod q,\qquad p<q^{3.99}.
\]
Consequently, there is an absolute constant $C>0$ such that
$P(a,q)\le Cq^{3.99}$ for all $q\ge2$ and $(a,q)=1$. In particular, Linnik's constant is at
most $3.99$.
\end{theorem}

We do not compute $q_0$ or $C$; the issue of effectivity is discussed at the end of
\cref{sec:join}. Except in the presence of an extremely close real zero, the proof finds a
prime in the smaller interval
\[
  q^{3.156342}<p<q^{3.99}.
\]
The lower endpoint comes from the support of the weight used to detect primes. When a
real zero is extremely close to $1$, Heath-Brown's theorem on Siegel zeros
\cite{HeathBrown1990siegel} instead gives $P(a,q)\le q^{3.5}$.

The argument combines published analytic estimates, new lemmas proved here, and exact
finite computations. The theorem requires all three parts. In particular, checking a
numerical certificate proves a bound for its specified configuration of zeros; the
analytic argument must also show that every modulus gives a configuration covered by
one of the certificates. We make this passage explicit in \cref{sec:casetree,sec:join}.

\subsection{Earlier work}
The admissible exponents in \cref{tab:history} record the main numerical developments.
Heath-Brown \cite{HeathBrown1992zero} obtained $5.5$, and Xylouris
\cite{Xylouris2009linnik,Xylouris2011least,Xylouris2011nullstellen,Xylouris2018linnik}
subsequently obtained $5.2$, $5.18$, and $5$. Before them, Pan \cite{Pan1957least,Pan1958least} gave
the first numerical values, and Chen, Jutila, Graham and Wang lowered them in a long series of
papers, which includes Graham's introduction of Selberg sieve weights into zero-density estimates
\cite{Graham1977thesis,Graham1981linnik}. Our proof uses the framework of these works. It rests on
the three principles that Heath-Brown isolates at the start of his paper: a zero-free region
with at most one exceptional zero
\cite{delaValleePoussin1896recherches,Gronwall1913series,Landau1918imaginar,Page1935number}, the
repulsion of other zeros by an exceptional zero, known as the Deuring--Heilbronn phenomenon
\cite{Deuring1933imaginare,Heilbronn1934class,Linnik1944deuring}, and a log-free zero-density
estimate \cite{Linnik1944least}. The last means that the density bound has no additional power of
$\log q$, a feature needed when studying zeros at distance $O(1/\log q)$ from $1$. Simpler proofs
of these principles, and of Linnik's theorem, were given by Rodosski\u\i{} \cite{Rodosskii1954least},
Tur\'an \cite{Turan1961density,Turan1984new}, Knapowski \cite{Knapowski1962linnik}, Fogels
\cite{Fogels1965zeros}, Gallagher \cite{Gallagher1970large}, Jutila \cite{Jutila1972density},
Motohashi \cite{Motohashi1975density,Motohashi1978primes} and Bombieri \cite{Bombieri1987grand},
among others; accounts are in the books \cite{Prachar1957primzahlverteilung,Montgomery1971topics,
IwaniecKowalski2004analytic,FriedlanderIwaniec2010opera}.

The conjectured scale is much smaller. Chowla \cite{Chowla1934least} conjectured
$P(a,q)\ll_\eps q^{1+\eps}$ for every $\eps>0$; the maximum over reduced classes is
conjecturally of order $\varphi(q)(\log q)^2$
\cite{Wagstaff1979greatest,Pomerance1980note,GranvillePomerance1990least,LiPrattShakan2017lower}.
The Generalized Riemann Hypothesis gives admissible exponents $2+\eps$; see
\cite{BachSorenson1996explicit,LamzouriLiSoundararajan2015conditional,CarneiroMilinovichQuesadaHerreraRamos2025fourier}
for explicit conditional bounds. Explicit \cite{BennettMartinOBryantRechnitzer2018explicit} and
uniform \cite{ThornerZaman2024refinements} unconditional estimates for primes in progressions are
also known. In the other direction, an exceptional (Siegel) zero helps: Heath-Brown
\cite{HeathBrown1990siegel} showed that a sufficiently strong exceptional zero forces
$P(a,q)\le q^{3+\delta}$, in an effective form that we use below, and $q^{2+\delta}$ ineffectively; see also
\cite{FriedlanderIwaniec2003exceptional,Iwaniec2006conversations,Zaman2016least}. Other proofs of
Linnik's theorem use pretentious methods \cite{Sachpazis2023pretentious}, sieve methods
\cite{FriedlanderIwaniec2023sifting}, or avoid $L$-functions
\cite{MatomakiMerikoskiTeravainen2024primes}, but give larger exponents. For moduli with special
multiplicative structure there are stronger zero-free regions
\cite{BarbanLinnikChudakov1964prime,Gallagher1972primes,Iwaniec1974zeros}, and much smaller
exponents are known \cite{Chang2014short}.

\begin{table}[ht]
\centering
\begin{tabular}{llll}
\toprule
$L$ & Year & Author(s) & Reference \\
\midrule
$10000$ & 1957 & Pan (announced) & \cite{Pan1957least} \\
$5448$ & 1958 & Pan & \cite{Pan1958least} \\
$777$ & 1965 & Chen & \cite{Chen1965least} \\
$630$ & 1971 & Jutila (reported by Tur\'an) & \cite{Turan1971recent} \\
$550$ & 1970 & Jutila & \cite{Jutila1970new} \\
$168$ & 1977 & Chen & \cite{Chen1977least} \\
$80$ & 1977 & Jutila & \cite{Jutila1977linnik} \\
$36$ & 1977 & Graham & \cite{Graham1977thesis} \\
$20$ & 1981 & Graham & \cite{Graham1981linnik} \\
$17$ & 1979 & Chen & \cite{Chen1979least} \\
$16$ & 1986 & Wang & \cite{Wang1986least} \\
$13.5$ & 1989 & Chen and Liu & \cite{ChenLiu1989III,ChenLiu1989IV} \\
$11.5$ & 1991 & Chen and Liu & \cite{ChenLiu1991V} \\
$8$ & 1991 & Wang & \cite{Wang1991least} \\
$5.5$ & 1992 & Heath-Brown & \cite{HeathBrown1992zero} \\
$5.2$ & 2009 & Xylouris & \cite{Xylouris2009linnik} \\
$5.18$ & 2011 & Xylouris & \cite{Xylouris2011least} \\
$5$ & 2011 & Xylouris & \cite{Xylouris2011nullstellen,Xylouris2018linnik} \\
\midrule
$3.99$ & 2026 & this paper & \\
\bottomrule
\end{tabular}
\caption{Admissible values of Linnik's constant, following the tables in
\cite{HeathBrown1992zero} and \cite{Xylouris2011least}, in order of the bound. Graham's value $20$
was submitted before Chen's $17$ appeared.}
\label{tab:history}
\end{table}

\subsection{From zeros to primes}\label{subsec:outline}
Write $\LL=\log q$ and express a zero of a Dirichlet $L$-function as
\[
  \rho=1-\frac{\lambda}{\LL}+i\frac{\mu}{\LL}.
\]
Thus small $\lambda$ means that the zero is close to the line $\sigma=1$; $\mu$ is its
height multiplied by $\log q$. We group each nonprincipal character with its complex
conjugate and call this a \emph{family}. The parameters $\lambda_1,\lambda_2,\lambda_3$
record the first zero of successive families in a specified rectangle near $1$.
The parameter $\lambda'$ records the next zero occurrence in the first family, after
removing the distinguished zero and its conjugate where appropriate. Precise definitions,
including multiplicities and ties, are given in \cref{sec:notation}.

A nonnegative weight $h$, supported in $[A,L]$ with $A=L-2T$ and $T=0.416829$, detects
primes between $q^A$ and $q^L$. Let $H$ be its Laplace transform. Xylouris's positivity
criterion \cite[(3.57)]{Xylouris2011nullstellen} gives
\[
  \sum_{p\equiv a\,(q)}\frac{\log p}{p}\,h\Bigl(\frac{\log p}{\LL}\Bigr)
  \ge \frac{\LL H(0)}{\varphi(q)}\bigl(1-W-\eta\bigr),
\]
where $\eta=10^{-6}$ and
\[
  W=\frac1{H(0)}\sum_{\chi\ne\chi_0}\ \sum_{\rho\in R_P}
       \bigl|H\bigl((1-\rho)\LL\bigr)\bigr|.
\]
Here $R_P$ is a fixed square in the normalized coordinates, and zeros are counted with
multiplicity. It is therefore enough to prove $W\le1-2\eta$. The contribution of an
individual zero decays exponentially with $\lambda$, so zeros close to $1$ require the
most accurate estimates.

We combine three kinds of information about these zeros. First, a \emph{per-character
bound} estimates the total contribution of one $L$-function from a lower bound for the
parameters of its zeros. Second, a \emph{far-density bound} limits a weighted sum over
characters, with weight depending on a selected zero of each character. Third,
\emph{near-density bounds} constrain the characters whose selected zeros lie closest to
$1$. Zero-location estimates for $\lambda_1$, $\lambda'$, $\lambda_2$, and $\lambda_3$ determine
which of these bounds are available in a given case.

\subsection{The new near-density estimate}\label{subsec:newlemma}
Heath-Brown's near-density estimate \cite[Lemma~12.1]{HeathBrown1992zero} bounds the number
$N(\lambda)$ of characters with a zero in $\sigma\ge1-\lambda/\LL$, $|t|\le1$. Its proof applies
the Cauchy--Schwarz inequality to a prime sum at the single point $\beta_1=1-\lambda_1/\LL$, and
compares the resulting Gram form of the characters with the diagonal. Because every character
is tested at the same point, a zero at distance $\lambda$ counts only through the single number
$N(\lambda)$, and the bound is weaker than the density estimate of Heath-Brown's \S11 as soon as
$\lambda\ge1.1$ and void soon after. That density estimate, proved with Selberg sieve weights
after Graham, has the form (11.4) of a sum over characters of a weight that depends on each
character's own distance $\lambda^{(k)}$, and it contains the count (11.5) of the characters with
$\lambda^{(k)}\le\lambda$. In his closing list of possible improvements Heath-Brown writes
\cite[\S16, item~9]{HeathBrown1992zero}:
\begin{quote}
It would be nice to have a weighted version of Lemma~12.1, in the way that (11.4) is a
weighted version of (11.5). Unfortunately, no neat way of achieving this seems available.
\end{quote}

The graded near-density lemma (\cref{thm:graded}) supplies such an estimate. We test each
character at its own \emph{anchor}, a real parameter determining the point at which a
smoothed prime sum is evaluated. Cauchy--Schwarz relates these prime sums to a Gram form
of the characters. We estimate its off-diagonal terms by the local explicit formula and,
where sieve weights are used, Burgess's bound; Graham's sieve asymptotic controls the
diagonal terms. Allowing the anchors to vary is the grading in the name of the lemma.

After normalization, the resulting constraint has the form
\[
  \Bigl(\sum_j a_jv_j\Bigr)_+^2
  \le \sum_j D_ja_j^2+d\Bigl(\sum_j a_j\Bigr)^2
  \qquad(a_j\ge0).
\]
Here $x_+=\max(x,0)$, the index $j$ runs over character--height entries, $v_j$ is a lower
bound for the response of an entry, and $D_j>0$ and $d>0$ bound its diagonal and common
correlation costs. Most entries correspond to one selected zero; some retain a second zero
of the same character, with an explicit correction for their correlation.

\Cref{prop:threshold} turns this inequality into a form suited to computation: there is a
number $\tau\in[0,\sqrt d]$ such that
\[
  \sum_j\frac{(v_j-\tau)_+^2}{D_j}\le1-\frac{\tau^2}{d}.
\]
A single constraint now distinguishes zeros at different distances from $1$.
The sieve weighting has precedents in Graham's work
\cite{Graham1977thesis,Graham1981linnik}, Heath-Brown's far-density argument, and the
sieve-weighted large-sieve inequalities of Motohashi \cite{Motohashi1975density}.

\emph{Where the new lemma enters.} In Heath-Brown's proof the near-density estimate enters at the very end, in the assembly of
\cite[\S15]{HeathBrown1992zero}. There the characters with a zero below a level $\Lambda$ are sorted into
bins $(\Lambda_{r+1},\Lambda_r]$, with $\Lambda_r=\Lambda-0.025r$; each character is charged the cost of the end of its
bin nearer to $s=1$; and the numbers of characters in the bins are controlled only through the
counts $N(\Lambda_r)$, bounded one level at a time by Lemma~12.1 (his Table~13). The far-density lemma
controls the characters beyond $\Lambda$. Xylouris's assembly \cite[\S6]{Xylouris2011nullstellen} has the same
structure, with a sharper form of Lemma~12.1 (his Lemma~5.3). A count $N(\lambda)$ treats all characters
with a zero below $\lambda$ alike, so it cannot tell a few zeros very close to $s=1$ from many zeros
further away. In our proof this place is taken by the \emph{near rows} of the leaf programs
described below. Each near row is one instance of \cref{thm:graded}: a single quadratic constraint
in which every bin of characters enters with its own feature, computed at its own anchor. The
weighted version of Lemma~12.1 that Heath-Brown asked for is thus used exactly where his count was
used, in every leaf program, and it is the main source of the improvement
(\cref{subsec:compare}).

\subsection{The finite case analysis}
We divide the middle range $0.1\le\lambda_1<1.5$ into cases determined by the type of the
first family and intervals for the first few zero parameters. There are four levels in
the computation, serving different purposes.
\begin{enumerate}[(1)]
  \item A \emph{root case} specifies a range of possible zero configurations.
    The $58$ parent rows taken from published zero-location tables lead to $4453$
    root cases, covering the middle range.
  \item A \emph{leaf} is a terminal case after further subdivision and valid
    zero-location exclusions. There are $4788$ leaves that require numerical bounds.
  \item A \emph{leaf program} groups characters into bins according to their selected
    zero parameters. Its nonnegative variables are the numbers of characters in these
    bins, together with weighted masses for the tails. The objective bounds $W$, and
    its constraints express the far-density bound, restrictions on the first families,
    and the near-density inequalities. For fixed thresholds these constraints are linear.
  \item A \emph{threshold box} bounds the auxiliary numbers $\tau$ in the near rows.
    On each box, a linear relaxation is either excluded or bounded by exact dual
    multipliers. Subdividing these boxes covers every possible choice of thresholds.
\end{enumerate}
The case tree partitions information about zeros; the threshold boxes partition the
auxiliary parameters used to bound them. Keeping these two subdivisions distinct is
essential to the coverage argument.

Each leaf, or each subcase of a leaf, has its own certificate, a binary tree of threshold
boxes. The $4788$ leaves give $5591$ such certificate trees, because some leaves are
subdivided further, and these trees contain $4{,}196{,}879$ threshold boxes. The largest certified value is
$0.9999998223\ldots<1$; it bounds the zero sum together with the required error allowance.
All certificate inequalities are checked in integer arithmetic. The tree was constructed
for an earlier computation at exponent $4.30$. Its subdivisions and zero-location
arguments are independent of the target exponent, so we retain them and certify new
leaf programs at $3.99$. \Cref{tab:scale} compares these numbers with the case analyses of
Heath-Brown and Xylouris, and \cref{tab:census} lists every level of the analysis.

The exterior ranges are shorter arguments. An exceptionally small $\lambda_1$ is covered
by Heath-Brown's theorem on Siegel zeros. The rest of $\lambda_1\le0.1$ is handled by his
quantitative Deuring--Heilbronn estimates and zero-location tables. A single further
certificate treats $\lambda_1\ge1.5$. If the relevant rectangle contains no zero, then
$W=0$ and the positivity criterion applies immediately.

\subsection{Comparison with Heath-Brown and Xylouris}\label{subsec:compare}

Our framework is Heath-Brown's \cite{HeathBrown1992zero} with Xylouris's refinements
\cite{Xylouris2011nullstellen,Xylouris2018linnik}. We use the following from their work:
\begin{itemize}
  \item the positivity criterion \cite[(3.57)]{Xylouris2011nullstellen};
  \item the per-character bound \cite[Lemma~3.10]{Xylouris2011nullstellen};
  \item the far-density lemma \cite[Lemma~5.1]{Xylouris2011nullstellen};
  \item the zero-free regions and Deuring--Heilbronn estimates, as printed in their tables;
  \item the local explicit formulas \cite[Lemmas~3.1, 3.2]{Xylouris2011nullstellen}, which are
    \cite[Lemmas~5.3, 5.2]{HeathBrown1992zero}.
\end{itemize}

\emph{The case analyses compared.} Both earlier proofs are computer-assisted, and both are
organized, like ours, as a case analysis on the first zeros; \cref{tab:scale} compares the three.
Heath-Brown's paper has twelve numerical tables of zero-location and density bounds. Its final
assembly \cite[\S15]{HeathBrown1992zero} treats $\lambda_1\ge0.348$ in fourteen intervals, mostly of length
$0.02$, each by a chain of inequalities evaluated numerically; the largest total it reports is
$0.9943$. Heath-Brown writes \cite[\S1]{HeathBrown1992zero} that ``the arguments rely heavily on
numerical calculations. These we do not reproduce in full, nor have we attempted any rigorous
analysis of the rounding and truncation errors in the computer algorithms employed.'' Xylouris's
dissertation adds eleven tables. Its final assembly \cite[\S6.4, pp.~86--87]{Xylouris2011nullstellen}
has $21$ cases: three intervals of $\lambda_1$ for a real zero of a real character, two for a complex
zero of a real character, and sixteen for a nonreal character. Each case is split further
according to the possible values of two zero-counting functions
\cite[(6.46)]{Xylouris2011nullstellen}, and the computations were done in Maple with ten digits
\cite[p.~14]{Xylouris2011nullstellen}; the largest value is $W<0.998$. In both, each case is closed by a
chain of inequalities in which the unknown zeros are estimated separately, one after another.

In our proof the same role is played by $4453$ root cases and $4788$ leaves, and each leaf is
closed by a linear program with an exact certificate instead of a floating-point evaluation. The
$4.2$ million boxes have no counterpart in the earlier work. They do not subdivide the
configurations of zeros. They subdivide the auxiliary thresholds of the near rows, and they arise
because the graded lemma is a quadratic constraint rather than a count. The largest certified
value, $0.99999982$ against $0.9943$ and $0.998$, shows how much finer the analysis must be at
$L=3.99$.

\begin{table}[tb]
\centering
\small
\setlength{\tabcolsep}{6pt}
\renewcommand{\arraystretch}{1.3}
\begin{tabular}{@{}>{\raggedright\arraybackslash}p{3.05cm}
                >{\centering\arraybackslash}p{3.25cm}
                >{\centering\arraybackslash}p{3.25cm}
                >{\columncolor{ourcolumn}\centering\arraybackslash}p{3.65cm}}
\toprule
 & \textbf{Heath-Brown} \cite{HeathBrown1992zero} & \textbf{Xylouris} \cite{Xylouris2011nullstellen} & \textbf{This paper}\\
 & $L=5.5$ \ (1992) & $L=5$ \ (2011) & $L=3.99$\\
\midrule
Near-density input & counts $N(\lambda)$ from Lemma~12.1 & counts $N(\lambda)$ from a sharper Lemma~12.1 &
  the weighted (graded) lemma, \cref{thm:graded}\\
Numerical tables & $12$ & $11$ more & $58$ rows from theirs, plus over $100$ new\\
Final case analysis & $14$ intervals of $\lambda_1$ & $21$ cases, each split by two zero counts &
  $4453$ roots, $4788$ leaves\\
How a case is closed & floating point, without rounding analysis & Maple, $10$-digit floating point &
  exact certificates on $4.2$~million boxes\\
Largest value of $W$ (must be below $1$) & $0.9943$ & $0.998$ & $0.99999982$\\
\bottomrule
\end{tabular}
\caption{Three proofs of an explicit Linnik constant compared, in the middle range of the first
zero ($\lambda_1\ge0.348$ for Heath-Brown and Xylouris, $0.1\le\lambda_1<1.5$ here). The last row is the largest
bound for the normalized zero sum $W$ over all cases of the final assembly: for Heath-Brown the
largest total he reports, for Xylouris his stated bound, and for this paper the largest certified
value, a bound for $W+2\eta$. The entry ``without rounding analysis'' paraphrases Heath-Brown's own
remark \cite[\S1]{HeathBrown1992zero}, quoted above.}
\label{tab:scale}
\end{table}

\emph{Where the gain comes from.} The gain from $5$ to $3.99$ has three sources:
\begin{itemize}
  \item the graded near-density lemma, which replaces the counts $N(\lambda)$ and applies to every
    range of $\lambda$ at once;
  \item the linear-programming combination, which lets all constraints act together on each
    case instead of chaining worst cases;
  \item a fine case analysis of the first zeros, possible because every case is certified by
    machine.
\end{itemize}
The last source alone does not go far. Xylouris estimates that his method, with much larger
computations (smaller intervals and finer grids), would give about $4.96$ \cite[p.~82]{Xylouris2018linnik}.
The earlier form of our method, with the linear programs and a fine case analysis but without the
graded lemma, reached $4.30$, using several ad hoc near-density models. The graded lemma replaces
almost all of them, and it is the main reason for the further gain to $3.99$.

\subsection{The certified exponent and numerical evidence}\label{subsec:why}
We state the result at $3.99$ because this is the exponent for which the complete cover
has been certified. The tightest cases have a nonreal first character and
$\lambda_1\in[0.64,0.80]$. In this range the available lower bounds for $\lambda_2$ are
comparatively small, which restricts the anchors in the near-density estimates.
Several of the corresponding leaf bounds are within $10^{-7}$ of $1$ after the error
allowance, and require fine subdivisions.

Exploratory floating-point computations suggest that a somewhat smaller exponent may
be accessible. One difficult leaf, with the reserved-family columns of
\cref{subsec:leafmodel}, has computed values $0.992$ at $L=3.97$ and $1.020$ at $L=3.95$.
These are values of particular relaxations, not certified global bounds or an obstruction
to other methods. Similarly, the positive numerical margins for the exterior ranges at
smaller exponents do not constitute a complete proof there. The verification reported in
this paper is at $L=3.99$.

\subsection{Verification and a guide to the proof}\label{subsec:checked}
The certificates have been checked by independent programs. Parts of the analytic
argument and the soundness of the checkers have also been formalized in Lean, with the
published analytic inputs stated as hypotheses. This formalization includes a theorem
that a valid certified leaf yields a prime when the modulus satisfies that leaf's
hypotheses about zeros. It does not include the full passage from arbitrary moduli to
leaves, or the exterior arguments. \Cref{sec:verification} gives the verification records
and separates these claims precisely.

\Cref{sec:notation,sec:inputs} give the notation and published inputs.
\Cref{sec:detection,sec:costs,sec:far} reduce prime detection to bounds for the zero sum,
and \cref{sec:location} supplies the zero-location estimates. The main new analytic
argument is in \cref{sec:near}; \cref{sec:rows} derives the constraints used in the
computation. \Cref{sec:lp} proves certificate soundness, while \cref{sec:casetree} shows
that the zero configurations are covered and satisfy the programs. The exterior cases
and the final choice of constants are treated in \cref{sec:exteriors,sec:join}.
A reader interested first in the new density estimate may read \cref{sec:notation,sec:inputs}
and then \cref{sec:near}, returning to the numerical construction afterwards.

\Cref{fig:proofmap} shows how the main results depend on one another, and which of them are
published inputs, new results, computations or formally verified statements. The terms of the
case analysis are introduced in \cref{subsec:vocabulary}, and the fixed constants and the error
budget in \cref{subsec:budget}. \Cref{tab:census} in \cref{sec:casetree} lists all levels of the
finite analysis, and \cref{tab:status} in \cref{sec:verification} records, for each component,
its source and how it was checked.

\begin{figure}[tbp]
\centering
\begin{tikzpicture}[
  x=1cm, y=1cm, font=\footnotesize,
  box/.style={rectangle, rounded corners=1.5pt, align=center, inner sep=3pt,
    minimum height=8.5mm, minimum width=17mm, line width=0.5pt},
  pub/.style={box, fill=black!12, draw=black!55},
  new/.style={box, fill=white, draw=black},
  comp/.style={box, fill=white, draw=black, dash pattern=on 2.2pt off 1.4pt},
  lean/.style={double, double distance=1pt, line width=0.45pt},
  dep/.style={-{Stealth[length=3.4pt,width=2.8pt]}, draw=black!75, line width=0.45pt},
  wire/.style={draw=black!75, line width=0.45pt},
]
% ---------------------------------------------------------------- theorem and regimes
\node[new] (main) at (8,0)
  {\textbf{Theorem~\ref{thm:main}}\\ $P(a,q)<q^{3.99}$ for $q\ge q_0$};
\node[new] (tiny) at (1.2,-1.8)
  {(R2) $\lambda_1<u_{\mathrm{exc}}$:\\ $P(a,q)\le q^{3.5}$\\ Prop.~\ref{prop:tiny}};
\node[new] (small) at (4.2,-1.8)
  {(R2)\\ $u_{\mathrm{exc}}\le\lambda_1<0.1$\\ Thm.~\ref{thm:small}};
\node[new] (r1) at (6.75,-1.8)
  {(R1) no zero\\ in $R(l)$,\\ so $W=0$};
\node[new, lean] (detect) at (9.7,-1.8)
  {$W\le1-2\eta$\\ gives a prime\\ Prop.~\ref{prop:detect}};
\node[pub] (crit) at (13.3,-1.8)
  {Xylouris's\\ criterion\\ Input~\ref{inp:criterion}};
\node[pub] (siegel) at (1.2,-3.6)
  {Siegel zeros\\ (Heath-Brown)\\ Input~\ref{inp:siegel}};
\node[pub] (hb) at (4.2,-3.6)
  {Heath-Brown's\\ lemmas\\ Inputs~\ref{inp:hb132}--\ref{inp:hb111}};
\node[comp] (marg) at (6.75,-3.6)
  {Interval\\ margins};
\node[new] (r3) at (9.7,-3.6)
  {(R3) $0.1\le\lambda_1<1.5$\\ Thm.~\ref{thm:cover}};
\node[new] (r4) at (12.6,-3.6)
  {(R4) $\lambda_1\ge1.5$\\ Thm.~\ref{thm:large}};
% ---------------------------------------------------------------- the middle range
\node[new] (src) at (1.0,-6.3)
  {4453 root\\ cases cover\\ Thm.~\ref{thm:sourcecover}};
\node[comp] (tree) at (3.3,-6.3)
  {Refinement\\ trees, \S\ref{subsec:refinement}\\ 4788 leaves};
\node[comp, lean] (certs) at (5.9,-6.3)
  {Exact\\ certificates:\\ $4{,}196{,}879$ boxes};
\node[new] (real) at (8.6,-6.3)
  {Realization:\\ feasible points\\ Prop.~\ref{prop:realization}};
\node[new, lean] (certthm) at (11.6,-6.3)
  {Certificate\\ soundness\\ Thm.~\ref{thm:certificate}};
\node[comp, lean] (largecert) at (14.4,-6.3)
  {One\\ certificate\\ (36 intervals)};
% ----
\node[new] (parent) at (1.0,-8.5)
  {58 parent\\ rows\\ Prop.~\ref{prop:parentrows}};
\node[new] (newloc) at (3.5,-8.5)
  {New location\\ bounds\\ Props.~\ref{prop:realrows}--\ref{prop:third}};
\node[new] (costs) at (6.3,-8.5)
  {Per-character\\ costs, \S\ref{sec:costs}};
\node[new] (farb) at (8.8,-8.5)
  {Far budget,\\ representatives\\ \S\ref{sec:far}};
\node[new] (rows) at (11.4,-8.5)
  {Near rows\\ hold\\ Prop.~\ref{prop:rowsvalid}};
% ----
\node[pub] (tables) at (1.3,-10.6)
  {Published\\ zero tables\\ Inputs~\ref{inp:rrtables}--\ref{inp:lambda3}};
\node[pub] (env) at (5.9,-10.6)
  {Envelope lemma\\ (Xylouris)\\ Input~\ref{inp:envelope}};
\node[pub] (farin) at (8.9,-10.6)
  {Far density\\ (Xylouris)\\ Input~\ref{inp:far}};
\node[new, lean] (zt) at (11.7,-10.6)
  {Threshold form\\ Prop.~\ref{prop:threshold},\\ Cor.~\ref{cor:zeroform}};
\node[new] (twotest) at (14.4,-10.6)
  {Two-test row\\ Thm.~\ref{thm:twotest}};
% ----
\node[new, lean] (response) at (10.1,-12.6)
  {Response lemma\\ Lem.~\ref{lem:response}};
\node[new, lean] (graded) at (13.5,-12.6)
  {\textbf{Graded}\\ \textbf{near-density lemma}\\ Thm.~\ref{thm:graded}};
\node[pub] (expl) at (10.3,-14.5)
  {Explicit formulas\\ Inputs~\ref{inp:X31}, \ref{inp:X32}};
\node[pub] (burgra) at (13.7,-14.5)
  {Burgess, Graham\\ Inputs~\ref{inp:burgess}, \ref{inp:graham}};
% ---------------------------------------------------------------- arrows
% both parts of (R2) give primes directly
\draw[wire] (tiny.north) |- (main.west);
\draw[dep] (small.north) |- (main.west);
\draw[dep] (detect.north) -- (detect.north |- main.south);
\draw[dep] (r1) -- (detect);
\draw[dep] (crit) -- (detect);
\draw[dep] (r3) -- (detect);
\draw[dep] (r4) -- (detect);
\draw[dep] (siegel) -- (tiny);
\draw[dep] (hb) -- (small);
\draw[dep] (marg) -- (small);
% the ingredients of the cover theorem, joined
\coordinate (bus) at ($(r3.south)+(0,-0.6)$);
\foreach \nd in {src,tree,certs,real}
  \draw[wire] (\nd.north) -- (\nd.north |- bus);
\draw[wire] ([xshift=-3mm]certthm.north) -- ([xshift=-3mm]certthm.north |- bus);
\draw[wire] (src.north |- bus) -- ([xshift=-3mm]certthm.north |- bus);
\draw[dep] (bus) -- (r3.south);
\draw[dep] ([xshift=4mm]certthm.north) -- ([xshift=4mm]certthm.north |- r4.south);
\draw[dep] (largecert) -- (r4);
% below
\draw[dep] (parent) -- (src);
\draw[dep] (newloc) -- (tree);
\draw[dep] (costs) -- (real);
\draw[dep] (farb) -- (real);
\draw[dep] (rows) -- (real);
\draw[dep] (tables) -- (parent);
\draw[dep] (tables) -- (newloc);
\draw[dep] (env) -- (costs);
\draw[dep] (farin) -- (farb);
\draw[dep] (zt) -- (rows);
\draw[dep] (twotest) -- (rows);
\draw[dep] (response) -- (zt);
\draw[dep] (graded) -- (zt);
\draw[dep] (expl) -- (response);
\draw[dep] (expl) -- (graded);
\draw[dep] (burgra) -- (graded);
% ---------------------------------------------------------------- legend
\node[pub, minimum width=7.5mm, minimum height=4.6mm] (k1) at (0.55,-12.3) {};
\node[new, minimum width=7.5mm, minimum height=4.6mm] (k2) at (0.55,-12.95) {};
\node[comp, minimum width=7.5mm, minimum height=4.6mm] (k3) at (0.55,-13.6) {};
\node[new, lean, minimum width=7.5mm, minimum height=4.6mm] (k4) at (0.55,-14.25) {};
\node[anchor=west] at ([xshift=1mm]k1.east) {published input};
\node[anchor=west] at ([xshift=1mm]k2.east) {result proved in this paper};
\node[anchor=west] at ([xshift=1mm]k3.east) {computation};
\node[anchor=west] at ([xshift=1mm]k4.east) {formally verified in Lean};
\end{tikzpicture}
\caption{Proof map for \cref{thm:main}. An arrow points from a result to a result whose
proof uses it; the joined lines into (R3) are the ingredients of \cref{thm:cover}. The
regimes (R1)--(R4) of \cref{sec:join} are defined by the first zero $\lambda_1$, and (R2)
splits at the constant $u_{\mathrm{exc}}$ of \cref{sec:exteriors}. A double border marks a
statement formally verified in Lean, with the published inputs as hypotheses, or, for the
certificates, a computation accepted by a checker proved sound in Lean. The formalization
also covers the far budget and parts of \cref{prop:rowsvalid,prop:realization}; the
zero-location results, the case tree, the costs, the two-test row, the exterior regimes and
the assembly are not formally verified (\cref{subsec:lean}). Some arrows are omitted, for
instance from the explicit formulas to \cref{sec:costs,sec:location} and into
\cref{thm:large}.}
\label{fig:proofmap}
\end{figure}

